% GENERATED FILE. Edit canonical lessons, then run npm run generate:books.
% canonical-source-hash: fnv1a64:09852d4ced16e8c8
% canonical-lessons: MW-C186-dhakkan, MW-C186-kudo, MW-C186-pidho, MW-C186-pato, MW-C186-palang

\chapter{Seats, Beds and Lids}
\label{ch:mw-seats-beds-and-lids}
\hlchaptermodality{186}

\begin{chapteropening}
\textbf{By the end of this chapter:} \emph{I can name the seats, the bed and the lids in a house.}

Complete the last lesson of chapter 186: I can name the seats, the bed and the lids in a house.
\end{chapteropening}

\section[\texorpdfstring{\mw{ढक्कण} (ḍhakkaṇ)}{ढक्कण (ḍhakkaṇ)}]{\mw{ढक्कण} (ḍhakkaṇ) — a lid}

\label{lesson:MW-C186-dhakkan}



\begin{quote}
Before the new one: say the Marwadi for asafoetida, then the Marwadi for fennel seed.
\end{quote}

\subsection*{You'll want to know: \mw{ढक्कण}}
\textbf{\mw{ढक्कण}} — \emph{ḍhakkaṇ} — \textquotedblleft{}a lid\textquotedblright{}.

\begin{quote}
\mw{हांडी} \mw{माथै} \mw{ढक्कण} \mw{है।}
\end{quote}

\begin{quote}
\emph{hā̃ḍī māthai ḍhakkaṇ hai.}
\end{quote}

\begin{quote}
There's a lid on the pot.
\end{quote}

A masculine noun.

\begin{cousinweb}
\textbf{\mw{ढक्कण}} (\emph{ḍhakkaṇ}) is built on a reconstructed \emph{*ḍhakk}, \textquotedblleft{}cover\textquotedblright{}, the word behind \textbf{\mw{ढांकणो}}; CDIAL lists Sindhi and Gujarati words for a lid under it, beside Marwari \emph{ḍhā̃kṇo} (CDIAL 5574).
\end{cousinweb}

\begin{grammarlens}[title={What you've built}]
One more word for this chapter.
\end{grammarlens}

\subsection*{Your turn}
\begin{itemize}
\raggedright
  \item \emph{Say it:} \emph{ḍhakkaṇ}
  \item \emph{Say it:} \emph{ḍhakkaṇ}, once more
  \item \emph{Say it:} \emph{hā̃ḍī māthai ḍhakkaṇ hai.} — There's a lid on the pot.
  \item \emph{Recall:} say the Marwadi for asafoetida, then the Marwadi for fennel seed, then say \emph{ḍhakkaṇ} again
\end{itemize}

\begin{tcolorbox}[breakable,colback=teal!4,colframe=teal!35!black,title={Before you move on}]
What does \mw{ढक्कण} mean? (\textquotedblleft{}A lid\textquotedblright{}.) Say it once more.
\end{tcolorbox}
\section[\texorpdfstring{\mw{कूंडो} (kū̃ḍo)}{कूंडो (kū̃ḍo)}]{\mw{कूंडो} (kū̃ḍo) — an earthen basin}

\label{lesson:MW-C186-kudo}



\begin{quote}
Before the new one: say the Marwadi for fennel seed, then the Marwadi for a lid.
\end{quote}

\subsection*{You'll want to know: \mw{कूंडो}}
\textbf{\mw{कूंडो}} — \emph{kū̃ḍo} — \textquotedblleft{}an earthen basin\textquotedblright{}.

\begin{quote}
\mw{कूंडो} \mw{भारी} \mw{है।}
\end{quote}

\begin{quote}
\emph{kū̃ḍo bhārī hai.}
\end{quote}

\begin{quote}
The earthen basin is heavy.
\end{quote}

A masculine noun.

\begin{cousinweb}
\textbf{\mw{कूंडो}} (\emph{kū̃ḍo}) is inherited from Sanskrit \emph{kuṇḍa}, \textquotedblleft{}a bowl, a waterpot\textquotedblright{}; CDIAL lists Hindi \emph{kū̃ḍā}, a clay pot that dough is kneaded in (CDIAL 3264). CDIAL gives no Marwari form in this entry.
\end{cousinweb}

\begin{grammarlens}[title={What you've built}]
One more word for this chapter.
\end{grammarlens}

\subsection*{Your turn}
\begin{itemize}
\raggedright
  \item \emph{Say it:} \emph{kū̃ḍo}
  \item \emph{Say it:} \emph{kū̃ḍo}, once more
  \item \emph{Say it:} \emph{kū̃ḍo bhārī hai.} — The earthen basin is heavy.
  \item \emph{Recall:} say the Marwadi for fennel seed, then the Marwadi for a lid, then say \emph{kū̃ḍo} again
\end{itemize}

\begin{tcolorbox}[breakable,colback=teal!4,colframe=teal!35!black,title={Before you move on}]
What does \mw{कूंडो} mean? (\textquotedblleft{}An earthen basin\textquotedblright{}.) Say it once more.
\end{tcolorbox}
\section[\texorpdfstring{\mw{पीढो} (pīḍho)}{पीढो (pīḍho)}]{\mw{पीढो} (pīḍho) — a low stool}

\label{lesson:MW-C186-pidho}



\begin{quote}
Before the new one: say the Marwadi for a lid, then the Marwadi for an earthen basin.
\end{quote}

\subsection*{You'll want to know: \mw{पीढो}}
\textbf{\mw{पीढो}} — \emph{pīḍho} — \textquotedblleft{}a low stool\textquotedblright{}.

\begin{quote}
\mw{रसोई} \mw{मांय} \mw{एक} \mw{पीढो} \mw{है।}
\end{quote}

\begin{quote}
\emph{rasoī mā̃y ek pīḍho hai.}
\end{quote}

\begin{quote}
There's a low stool in the kitchen.
\end{quote}

A masculine noun.

\begin{cousinweb}
\textbf{\mw{पीढो}} (\emph{pīḍho}) is inherited from Sanskrit \emph{pīṭha}, \textquotedblleft{}a stool, a bench\textquotedblright{}; CDIAL lists Hindi \emph{pīṛhā}, \textquotedblleft{}a large square stool\textquotedblright{} (CDIAL 8222). CDIAL gives no Marwari form in this entry.
\end{cousinweb}

\begin{grammarlens}[title={What you've built}]
One more word for this chapter.
\end{grammarlens}

\subsection*{Your turn}
\begin{itemize}
\raggedright
  \item \emph{Say it:} \emph{pīḍho}
  \item \emph{Say it:} \emph{pīḍho}, once more
  \item \emph{Say it:} \emph{rasoī mā̃y ek pīḍho hai.} — There's a low stool in the kitchen.
  \item \emph{Recall:} say the Marwadi for a lid, then the Marwadi for an earthen basin, then say \emph{pīḍho} again
\end{itemize}

\begin{tcolorbox}[breakable,colback=teal!4,colframe=teal!35!black,title={Before you move on}]
What does \mw{पीढो} mean? (\textquotedblleft{}A low stool\textquotedblright{}.) Say it once more.
\end{tcolorbox}
\section[\texorpdfstring{\mw{पाटो} (pāṭo)}{पाटो (pāṭo)}]{\mw{पाटो} (pāṭo) — a plank seat}

\label{lesson:MW-C186-pato}



\begin{quote}
Before the new one: say the Marwadi for an earthen basin, then the Marwadi for a low stool.
\end{quote}

\subsection*{You'll want to know: \mw{पाटो}}
\textbf{\mw{पाटो}} — \emph{pāṭo} — \textquotedblleft{}a plank seat\textquotedblright{}.

\begin{quote}
\mw{पाटो} \mw{नीचो} \mw{है।}
\end{quote}

\begin{quote}
\emph{pāṭo nīcho hai.}
\end{quote}

\begin{quote}
The plank seat is low.
\end{quote}

A masculine noun.

\begin{cousinweb}
\textbf{\mw{पाटो}} (\emph{pāṭo}) is inherited from Sanskrit \emph{paṭṭa}, \textquotedblleft{}a slab, a tablet\textquotedblright{}; CDIAL lists Old Marwari \emph{pāṭī}, \textquotedblleft{}a plank\textquotedblright{}, and a Hindi board to sit on while eating (CDIAL 7699).
\end{cousinweb}

\begin{grammarlens}[title={What you've built}]
One more word for this chapter.
\end{grammarlens}

\subsection*{Your turn}
\begin{itemize}
\raggedright
  \item \emph{Say it:} \emph{pāṭo}
  \item \emph{Say it:} \emph{pāṭo}, once more
  \item \emph{Say it:} \emph{pāṭo nīcho hai.} — The plank seat is low.
  \item \emph{Recall:} say the Marwadi for an earthen basin, then the Marwadi for a low stool, then say \emph{pāṭo} again
\end{itemize}

\begin{tcolorbox}[breakable,colback=teal!4,colframe=teal!35!black,title={Before you move on}]
What does \mw{पाटो} mean? (\textquotedblleft{}A plank seat\textquotedblright{}.) Say it once more.
\end{tcolorbox}
\section[\texorpdfstring{\mw{पलंग} (palaṅg)}{पलंग (palaṅg)}]{\mw{पलंग} (palaṅg) — a bedstead}

\label{lesson:MW-C186-palang}



\begin{quote}
Before the new one: say the Marwadi for a low stool, then the Marwadi for a plank seat.
\end{quote}

\subsection*{You'll want to know: \mw{पलंग}}
\textbf{\mw{पलंग}} — \emph{palaṅg} — \textquotedblleft{}a bedstead\textquotedblright{}.

\begin{quote}
\mw{म्हारो} \mw{पलंग} \mw{नुवो} \mw{है।}
\end{quote}

\begin{quote}
\emph{mhāro palaṅg nuvo hai.}
\end{quote}

\begin{quote}
My bed is new.
\end{quote}

A masculine noun.

\begin{cousinweb}
\textbf{\mw{पलंग}} (\emph{palaṅg}) is inherited from Sanskrit \emph{palyaṅka}, \textquotedblleft{}a bed\textquotedblright{}; CDIAL lists an Old Marwari word for a bed under it (CDIAL 7964). It is a framed bed; \textbf{\mw{खाट}} is the string bed.
\end{cousinweb}

\begin{grammarlens}[title={What you've built}]
One more word for this chapter.
\end{grammarlens}

\subsection*{Your turn}
\begin{itemize}
\raggedright
  \item \emph{Say it:} \emph{palaṅg}
  \item \emph{Say it:} \emph{palaṅg}, once more
  \item \emph{Say it:} \emph{mhāro palaṅg nuvo hai.} — My bed is new.
  \item \emph{Recall:} say the Marwadi for a low stool, then the Marwadi for a plank seat, then say \emph{palaṅg} again
\end{itemize}

\begin{tcolorbox}[breakable,colback=teal!4,colframe=teal!35!black,title={Before you move on}]
What does \mw{पलंग} mean? (\textquotedblleft{}A bedstead\textquotedblright{}.) Say it once more.
\end{tcolorbox}
